% =====================================================================
%  TEMPLATE — /note  and  /note-with-qr     (book-style lecture note)
%  /note          -> qr=false        /note-with-qr -> qr=true
%  Replace every <<...>>. Delete blocks you do not need. Build with:
%     bash <skill>/scripts/build.sh note.tex transcript.tsv
% =====================================================================
\documentclass[11pt]{article}
\usepackage{inote}
\inotesetup{mode=note, qr=true,
  video=<<11-char YouTube ID>>,                 % from: transcript.py id <url>
  brand={<<brand / channel name or leave empty>>},
  subject={<<HSC Physics 1st Paper · Chapter-02>>},
  chapter={<<ভেক্টর>>}, chapter-en={<<Vector>>}, chapter-no={<<দ্বিতীয় অধ্যায়>>}}
\begin{document}
\chapterbanner                    % full-lecture QR sits in the right chevron

% ---------- 1. LECTURE MAP (one row per Type, in lecture order) ----------
\sectionbar{Lecture Map — টপিক অনুযায়ী ভাগ}
\begin{tblr}{colspec={X[2.2,l,m] Q[c,m] Q[c,m] X[2.6,l,m] Q[c,m]},
  hlines={gridgray,.5pt},vlines={gridgray,.5pt},
  row{1}={bg=headgray,font=\labelfont\bfseries,halign=c},rowsep=2.5pt}
টাইপের নাম & গুরুত্ব & ভিডিওতে & যা যা আছে & পৃষ্ঠা\\
\textbf{Type-01:} <<name>> & \stars{5} & <<03:10>> & <<what is inside>> & <<২>>\\
\end{tblr}
% (Board-analysis table ONLY if the user's source contains real board data.)

% ---------- 2. FORMULA SHEET: every formula of the lecture, ONCE ----------
\fancyheading{গুরুত্বপূর্ণ সূত্রাবলি ও বিশ্লেষণ}
\begin{formulagrid}[0.52]
\sqitem{<<formula group title>>:}
$<<formula>>$\par
\hspace*{3mm}<<symbol meanings>>
&
<<diagram (tikzpicture) OR where-clause / analysis bullets>>\\
\end{formulagrid}

\begin{unitbox}[<<একক রূপান্তর / দ্রুত মনে রাখো>>]
\item <<one fact per item>>
\end{unitbox}

% ---------- 3. TYPE-WISE: concept first, then the maths of that type ----------
\begin{multicols}{2}
\begin{typebox}{01}{<<Type name>>}
\withqr{<<sec where teacher STARTS explaining this concept>>}{%
\textbf{কনসেপ্ট:} <<explanation in simple Bangla, built from the teacher's words>>}
\sqitem{<<key formula name>>:} $<<formula>>$
\begin{roundbox}\textbullet\ <<key takeaway>>\end{roundbox}
\end{typebox}

\begin{tipbox}[<<sec>>]
<<a tip the teacher SAID that is NOT on the slide>>
\end{tipbox}

\cq{১}[<<source: ভিডিও ম্যাথ / বোর্ড + সাল if known>>]{<<question stem>>}[<<sec where this math starts>>]
\solhead
\workgiven[0.62]{%
$<<step 1>>$\\
$<<step 2>>$\\
$<<answer>>$ \Ans}
{\given{$<<data 1>>$\\ $<<data 2>>$}}

% CQ with parts:
% \cq{২}[...]{<<stem>>}[<<sec>>]
% \qpart{গ}{<<part c>>}  \qpart{ঘ}{<<part d>>}
% \solhead \anslabel{গ} ... \anslabel{ঘ} ...
\end{multicols}
\end{document}
